%% ma: L11-B03-0014
%% lop: 11
%% chuong: 1
%% bai: 3
%% dang: 11.3.6
%% loai: TL
%% muc: VD
%% dap_an: "-"
%% nguon: "de-cuong-giua-ky-1-toan-11-nam-2026-2027-truong-thpt-hoang-van-thu-ha-noi-2-8.pdf (M08, Đề cương giữa kì 1 Toán 11 THPT Hoàng Văn Thụ, Hà Nội), Phần 4 Bài 3"
%% kien_thuc: "Bài 1 (giá trị lượng giác của góc lượng giác), Bài 3 (hàm số lượng giác), hàm số bậc hai trên một đoạn (lớp 10)"
%% trang_thai: cho-duyet
%% ly_do: "Dạng mới đề xuất 11.3.6"
\begin{ex}
Tìm giá trị lớn nhất, giá trị nhỏ nhất của các hàm số sau:

a) $y=3\sin x-13$;\quad b) $y=\sin^2x+\cos x-2026$;\quad c) $y=5-2\sin^22x$.
\loigiai{a) $-1\le\sin x\le1\Rightarrow-16\le y\le-10$. Giá trị lớn nhất $-10$ (khi $x=\dfrac\pi2+k2\pi$), giá trị nhỏ nhất $-16$ (khi $x=-\dfrac\pi2+k2\pi$).

b) Đặt $t=\cos x\in[-1;1]$: $y=-t^2+t-2025$. Trên $[-1;1]$, hàm bậc hai đạt giá trị lớn nhất tại $t=\dfrac12$: $y=-\dfrac{8099}{4}$; giá trị nhỏ nhất tại $t=-1$: $y=-2027$. Vậy $\max y=-\dfrac{8099}4$ (khi $\cos x=\dfrac12$), $\min y=-2027$ (khi $\cos x=-1$).

c) $0\le\sin^22x\le1\Rightarrow3\le y\le5$. Giá trị lớn nhất $5$ (khi $\sin2x=0$), giá trị nhỏ nhất $3$ (khi $\sin^22x=1$).}
\end{ex}
